% === Related Work ===
\section{Related Work}
\label{sec:related}

\noindent\textbf{Incident management and AIOps.} Modern cloud and microservice
operations generate incident volumes that overwhelm manual on-call triage,
motivating a large body of AIOps research surveyed by He et
al.~\cite{he2021survey}. Empirical studies characterize the human cost of this
process: Ghosh et al.~\cite{ghosh2022incidents} analyze a large-scale cloud
service and document the long detection-to-mitigation pipeline that motivates
automation, while Chen et al.~\cite{chen2019triage} frame \emph{continuous
incident triage} as a learning problem for routing incidents to the correct team.
\sys{} targets the same triage bottleneck but reframes it as retrieval against a
memory of resolved incidents rather than classification.

\noindent\textbf{LLMs for incident diagnosis.} Recent work applies large language
models to root-causing and mitigation. Ahmed et al.~\cite{ahmed2023rootcause}
perform the first large-scale study of LLMs recommending root-cause and mitigation
steps; RCACopilot~\cite{chen2024rcacopilot} couples diagnostic-data collection
with LLM reasoning and matches incoming incidents to handlers by alert type; and
D-Bot~\cite{zhou2024dbot} diagnoses database faults with LLM agents.
AIOpsLab~\cite{chen2025aiopslab} provides a framework for evaluating such agents
in autonomous-cloud settings. \sys{} differs by introducing a capability-gated
controller that runs only the reasoning steps the available evidence supports,
degrading gracefully from full telemetry to text-only inputs, and by suppressing
duplicate pages within an incident.

\noindent\textbf{Log analysis and anomaly detection.} A complementary line of
work mines telemetry directly. Drain~\cite{he2017drain} parses unstructured logs
into templates; DeepLog~\cite{du2017deeplog}, LogAnomaly~\cite{meng2019loganomaly},
and LogBERT~\cite{guo2021logbert} detect sequential and semantic anomalies; and
benchmark resources such as the log-parsing toolkit~\cite{zhu2019logparsing}, the
foundational anomaly-detection experience report~\cite{he2016experience}, and the
Loghub collection~\cite{zhu2023loghub} standardize evaluation. \sys{} consumes
parsed telemetry as one capability tier but does not depend on it, retaining
text-only operation when traces and metrics are unavailable.

\noindent\textbf{Retrieval and retrieval-augmented reasoning.} \sys{}'s retriever
fuses dense, sparse, and graph signals. Dense bi-encoders build on
Sentence-BERT~\cite{reimers2019sbert} and its backbones
MiniLM~\cite{wang2020minilm} and MPNet~\cite{song2020mpnet}, with general-purpose
embedding families E5~\cite{wang2022e5} and BGE/C-Pack~\cite{xiao2024cpack};
sparse retrieval uses BM25~\cite{robertson2009bm25} and the learned-sparse
SPLADE~\cite{formal2021splade}; cross-encoder re-ranking follows Nogueira and
Cho~\cite{nogueira2019passage}; and we combine ranked lists with reciprocal rank
fusion~\cite{cormack2009rrf}. Dense Passage Retrieval~\cite{karpukhin2020dpr} and
Retrieval-Augmented Generation~\cite{lewis2020rag} established the
retrieve-then-reason paradigm, and ReAct~\cite{yao2023react} interleaves reasoning
with tool actions; \sys{}'s evidence-gathering loop and its underlying
LLM~\cite{qwen2025} inherit this lineage.

\noindent\textbf{Knowledge graphs and duplicate-incident retrieval.} Building
knowledge graphs from incidents and tickets supports structured retrieval;
SoftNER~\cite{shetty2022softner} mines such graphs from cloud incidents, informing
\sys{}'s knowledge-graph retriever. \sys{}'s ``retrieve the most similar past
incident'' framing also descends directly from duplicate and similar bug-report
retrieval: Runeson et al.~\cite{runeson2007duplicate} pioneered NLP-based
duplicate-defect detection, Sun et al.~\cite{sun2010duplicate} introduced a
discriminative retrieval model, and recent deep-learning approaches such as
dual-channel CNNs~\cite{he2020dccnn} improve matching. \sys{} extends this idea
from bug reports to operational incidents and adds telemetry-aware,
capability-gated retrieval.

\noindent\textbf{Benchmark systems and data.} Our two collected
microservice-telemetry datasets are built on Google's Online
Boutique~\cite{googleboutique} and the OpenTelemetry Demo~\cite{oteldemo} under
controlled fault injection with Chaos Mesh~\cite{chaosmesh}. Our real-world corpus
is drawn from the World of Logs collection~\cite{wang2026worldoflogs}, a public
archive of logs and issue reports mined from online software-engineering sources;
other large public Jira datasets, such as that of Montgomery et
al.~\cite{montgomery2022jira}, offer complementary issue data.
